%% Chapter 1 - Introduction

\chapter{Introduction}
\label{ch:introduction}

Dental panoramic radiography, also referred to as Orthopantomography (OPG),
provides a single two-dimensional view of the full dentition, alveolar bone,
maxilla, mandible, and surrounding anatomical structures. Because one image
captures most of the clinically relevant oral region, OPGs are routinely used
for screening dental caries, periodontal bone loss, impacted teeth, and other
oral pathologies. Their practical value is high, but so is the difficulty of
interpretation: weak local contrast, overlapping structures, projection
distortion, and subtle early lesions all increase the risk of missed or
inconsistent findings.

The present project aims to build an intelligent dental assisting system that
uses deep learning to move from a full OPG to a structured tooth-wise report.
The target system is tooth-centric. Instead of classifying the full radiograph
as a single unit, the OPG is first decomposed into tooth-level regions and each
tooth is then analyzed through task-specific models. This design is clinically
motivated: dentists reason per tooth, and an assistive system is only useful if
it can localize its findings and map them to a standard numbering scheme.

The broader survey work conducted alongside this project also reinforced a
system-level conclusion that directly shaped Phase 2: current dental AI
literature is strong on isolated subtasks, but weak on complete OPG-to-report
workflows. This is why the present report emphasizes modular implementation and
honest integration status instead of claiming a fully finished clinical system.

At the project level, the intended end-to-end workflow is: full OPG input,
tooth detection, FDI numbering, caries analysis on tight tooth crops,
periodontal analysis on expanded bone-aware tooth crops, and final report
generation. However, this report is restricted to work completed up to the
current Phase 2 milestone. The major completed contributions are the
dataset-engineering and model-development steps for the caries and periodontal
branches. On the caries side, a collaborator has built a two-stage ResNet18
cascade on DENTEX tooth crops. On the periodontal side, the work completed in
this repository includes Tufts annotation exploration, correction of critical
data parsing bugs, tooth-level crop generation, overlap-based label assignment,
binary bone-loss classification, overlap-based severity classification, and an
exploratory CEJ-based estimation module.

One of the most important outcomes of this phase is that the periodontal branch
is already supported by reproducible code and measured results. The bone-loss
binary classifier trained on resampled Tufts crops achieved a validation
macro-F1 of 0.9233, while the overlap-based severity model achieved a macro-F1
of 0.7208. The CEJ-based heuristic branch was investigated because it is closer
to the clinical standard, but the current OPG-based formulation remains
insufficiently reliable and is therefore treated as future work rather than the
primary reporting pathway.

This report is organized as follows. Chapter~\ref{ch:architecture} describes
the target system architecture and clarifies which parts are already
implemented. Chapter~\ref{ch:methodology} documents the datasets, preprocessing
choices, model configurations, and training design used so far.
Chapter~\ref{ch:implementation} presents the implemented modules in the current
repository and notebooks. Chapter~\ref{ch:results} summarizes the quantitative
and qualitative findings supported by the saved artifacts. Finally,
Chapter~\ref{ch:conclusion} concludes the report and outlines the most
immediate next steps toward full pipeline integration.
